\section{Identification of the physical transition}
\label{sec:transition}

The construction in Proposition~\ref{flow:tuned} selects a parameter
$t_*=t_*(D)$, and Propositions~\ref{read:moments} and~\ref{read:arm}
give correlation moments and a wired one-arm estimate there.  We now show
that $t_*=0$ and that this value is the magnetization threshold.  The
first assertion uses partition functions on tori and a product-space
influence inequality.  The second uses the high moment in
Proposition~\ref{read:arm} to control a decision-tree exploration in a fixed
environment.

For a finite graph $G=(V,E)$, put $v_e=e^{2K_e}-1$ and write
\begin{equation}\label{crit:fk-law}
 \phi_{G,t,\omega}(A)
 =\frac{1}{\mathcal Z_{G,t,\omega}}
       2^{k(A)}\prod_{e\in A}v_e,
 \qquad A\subset E.
\end{equation}
Here $k(A)$ counts components of the spanning subgraph $(V,A)$; a
specified boundary wiring is imposed before counting components.
Superscripts $\mathrm f$ and $\mathrm w$ indicate free and wired
boundaries.  All these probabilities are normalized at fixed $\omega$.
The Edwards--Sokal coupling~\cite{EdwardsSokal1988} identifies the
wired connection probability to the boundary with the corresponding
plus-boundary spin expectation.  In infinite volume it gives
\begin{equation}\label{crit:magnetization}
 \langle\sigma_0\rangle^+_{t,\omega}
       =\phi^{\mathrm w}_{t,\omega}(0\longleftrightarrow\infty).
\end{equation}
Finite-energy bounds, positive association, and comparison in the
couplings hold uniformly on any fixed compact $(t,D)$ window considered
below.  In particular, for some $\eta>0$, every one-edge conditional
open probability, with any other edges prescribed, lies in
$[\eta,1-\eta]$.  This follows directly from
\eqref{crit:fk-law}: that probability is either $v_e/(1+v_e)$ or
$v_e/(2+v_e)$.

\subsection{Torus duality and twisted partition functions}

Let $\mathbb T_\ell=(\Z/\ell\Z)^2$ have its usual square cellulation.
For $A\subset E(\mathbb T_\ell)$, let
$H(A)\subset H_1(\mathbb T_\ell;\mathbb F_2)=\mathbb F_2^2$ be the
image of its cycle space, and put $r(A)=\dim H(A)$.  The complementary
dual configuration $A^*$ contains the dual of each edge outside $A$.
For each nonzero character
$\chi:\mathbb F_2^2\longrightarrow\mathbb F_2$, define the increasing,
translation-invariant event
\[
 \mathcal H_\chi=\{A:\chi|_{H(A)}\ne0\}.
\]
Thus $\mathcal H_\chi$ asks for an open cycle detected by $\chi$.

\begin{lemma}[Torus duality]\label{crit:torus-duality}
Under the intersection pairing,
\begin{equation}\label{crit:homology-annihilator}
 H(A^*)=H(A)^\perp,\qquad r(A^*)=2-r(A).
\end{equation}
If $v_{e^*}^*=2/v_e$, then complement duality carries the torus FK law
at $v$ to a measure whose density with respect to the torus FK law at
$v^*$ lies between $1/4$ and $4$.
\end{lemma}

\begin{proof}
The cycle space of $(V,A)$ has dimension $|A|-|V|+k(A)$.  A cellular
boundary is supported on $A$ precisely when its face coefficients are
constant on each complementary dual component.  The space of such
boundaries therefore has dimension $k(A^*)-1$: its face coefficients
have $k(A^*)$ degrees of freedom, and adding the constant two-chain
does not change the boundary.  Consequently
\begin{equation}\label{crit:euler}
 |A|-|V|+k(A)=k(A^*)-1+r(A).
\end{equation}
Apply the same identity on the dual cellulation and add the two
identities.  The torus Euler formula $|V|-|E|+|F|=0$ gives
$r(A)+r(A^*)=2$.  Primal and complementary dual cycles have zero
intersection, since a primal edge and its transverse dual cannot both
be present.  This inclusion, together with the dimension identity,
proves \eqref{crit:homology-annihilator}.

The ratio of the two unnormalized weights is
\begin{align*}
 \frac{2^{k(A)}\prod_{e\in A}v_e}
 {2^{k(A^*)}\prod_{e\notin A}(2/v_e)}
 &=\Big(\prod_{e\in E}v_e\Big)
            2^{k(A)-k(A^*)-|E|+|A|}\\
 &=\Big(\prod_{e\in E}v_e\Big)
            2^{|V|-|E|-1}\,2^{r(A)},
\end{align*}
where the last equality uses \eqref{crit:euler}.  The first factor is
independent of $A$, and $1\le2^{r(A)}\le4$.  Normalization proves the
claimed comparison.
\end{proof}

Choose signs $s_e^\chi\in\{-1,1\}$ representing the cocycle $\chi$,
and let $Z^\chi(K)$ be the Ising partition function with interactions
$K_es_e^\chi\sigma_x\sigma_y$.  Write $Z^0(K)$ for the ordinary
partition function.  Expanding
\[
 e^{K_es_e^\chi\sigma_x\sigma_y}
 =e^{-K_e}\bigl(1+v_e\one_{\{\sigma_x=s_e^\chi\sigma_y\}}\bigr)
\]
shows that a component of $A$ contributes two compatible spin
assignments exactly when the product of seam signs around each of its
cycles is $1$.  An incompatible component contributes zero.  Hence,
for every environment, not only after averaging,
\begin{equation}\label{crit:twist-identity}
 \frac{Z^\chi(K)}{Z^0(K)}
          =\phi_{\mathbb T_\ell,t,\omega}(\mathcal H_\chi^c).
\end{equation}

At the critical reference couplings there is $a_0>0$, independent of
$\ell$, such that
\begin{equation}\label{crit:reference-homology}
 a_0\le\phi^{\mathrm{ref}}_{\mathbb T_\ell}(\mathcal H_\chi)
                    \le1-a_0
 \qquad(\chi\ne0).
\end{equation}
For completeness, choose a coordinate generator on which $\chi$ is
nonzero.  In a narrow toric strip along that generator, a bounded
number of overlapping fixed-aspect rectangle crossings, with transverse
crossings in their overlaps, produces an open generator cycle.
The arbitrary-boundary crossing bounds of
\cite[Theorem~1.1]{CDH2016}, applied in the rectangular charts, and FK
positive association give a uniform lower bound.  Make both generator
cycles in the dual configuration in the same way.  Their homology has
rank two, so \eqref{crit:homology-annihilator} forces $H(A)=0$.
Lemma~\ref{crit:torus-duality} compares the complemented law with the
critical dual FK law by an absolute factor.  This proves the upper
bound in \eqref{crit:reference-homology}.  Only finitely many crossings
are used, so their lower bounds do not deteriorate with $\ell$.

\begin{lemma}[Transfer of torus nondegeneracy]\label{crit:torus-transfer}
Fix a sufficiently large integer $Q$.  After choosing the smallness
constants of the construction sufficiently small, there is $a_1>0$
such that, for $\ell_J=Qr_J$ and every nonzero $\chi$,
\begin{equation}\label{crit:tuned-homology}
 a_1\le \E\phi_{\mathbb T_{\ell_J},t_*,\omega}(\mathcal H_\chi)
                 \le1-a_1
\end{equation}
for all sufficiently large $J$.
\end{lemma}

\begin{proof}
Apply the exact finite-volume map of Proposition~\ref{block:finite}
to the untwisted model and to the $\chi$-twisted model, using the
corresponding critical reference partition function in each case and
the same decision seeds.  A twist is removed on a contractible chart
by a spin gauge transformation.  Every background function is even,
so the ambiguity of a simultaneous sign change on the chart does not
alter it.  Conditional reference averaging, the local decisions, and
plane centering in this gauge therefore give identical chart
coefficients and identical saved scalar increments in the two runs.

Here is why the errors outside charts remain uniform in $J$.  Choose
$Q$ larger than the fixed halos and chart margins in the construction.
A principal output uses only a short parent footprint, by
Lemma~\ref{block:carrier}.  That footprint lifts to the plane at
every scale up to $J$, and its environment and seeds have the planar
law.  A nonchart footprint has length at least $cm$ when the torus has
side $m$ at the scale carrying that footprint.  Changing between input
and output cells changes this constant by the fixed factor $L$ only.
It cannot satisfy the short-footprint rule
for a principal output.  This also proves that a seam crossing a
short footprint causes no exception: it is removed by the local
gauge, whereas a genuinely nonchart input remains in the tail.

Keep the scalar supports until applying the enhanced tree weights in
Lemma~\ref{block:tree} and Corollary~\ref{block:diameter}.
Summing the possible nonchart footprints,
including the neighborhoods used by the decisions, bounds the
uncancelled scalar increment at this step by
\[
 C_L\delta m^2e^{-cm}.
\]
There are $O(m^2)$ possible roots; the additional length weight pays
$e^{-cm}$ after the animal sum.  The side lengths $m$ decrease
geometrically and remain at least $Q$, so their total is bounded by
$C_L\delta Q^2e^{-c'Q}$, independently of $J$.  The final tail functions
have total supremum norm at most $C\delta Q^2$.  These estimates use
the finite-volume weighted bounds before forgetting supports, and
apply also to normalized conditional operations on nonchart terms.

To see precisely how these errors enter the desired ratio, let
$c_J^\chi,H_J^\chi$ denote the accumulated scalar and terminal
background for a twist, and use superscript $0$ for the untwisted
run.  Exact partition preservation gives
\begin{equation}\label{crit:terminal-ratio}
 \frac{Z^\chi(K)}{Z^0(K)}
 =\frac{Z^\chi(K^{\mathrm{ref}})}{Z^0(K^{\mathrm{ref}})}
       e^{c_J^\chi-c_J^0}
       \frac{\langle e^{H_J^\chi}\rangle_{J,\mathrm{ref},\chi}}
            {\langle e^{H_J^0}\rangle_{J,\mathrm{ref},0}}.
\end{equation}
The brackets are normalized terminal reference expectations in the
respective exact partition identities.  Chart scalar increments
cancel in $c_J^\chi-c_J^0$, leaving the nonchart bound above.
At the final scale there are only $O(Q^2)$ principal slots.  Their
laws agree in charts with the planar laws, so the high-moment bounds
along the tuned flow imply that their total supremum norm converges
to zero in probability as $J\to\infty$.  Denote this total, in the two
runs together, by $P_J$.  Exact partition preservation and
$|\log\langle e^H\rangle|\le\norm{H}_\infty$ now give
\begin{equation}\label{crit:log-ratio-transfer}
 \left|\log\frac{Z^\chi(K)/Z^0(K)}
                         {Z^\chi(K^{\mathrm{ref}})/Z^0(K^{\mathrm{ref}})}
       \right|
 \le P_J+C\delta Q^2+C_L\delta Q^2e^{-c'Q},
 \qquad P_J\xrightarrow{\Prob}0.
\end{equation}
The reference ratios here are on the same microscopic torus; exact
blocking has not replaced them by a limiting torus ratio.
First fix $Q$, then choose $\delta$ and the permitted disorder so that
the deterministic errors are smaller than a fixed tolerance depending
on $a_0$, and finally let $J$ grow.  By
\eqref{crit:twist-identity} and \eqref{crit:reference-homology}, on an
event of probability tending to one all three tuned character
probabilities lie in a common compact subinterval of $(0,1)$.
Their values always lie in $[0,1]$, so averaging proves
\eqref{crit:tuned-homology}.
\end{proof}

\subsection{Sharp thresholds after averaging the environment}

We next prove the threshold estimate needed for the torus events.
The environment is independent across bonds, but the FK edges are
not.  A monotone sequential representation lets us apply the product
influence theorem while keeping the useful derivative as a quenched
covariance.

\begin{lemma}[Torus threshold inequality]\label{crit:threshold}
Fix a compact parameter window containing $0$ and $t_*(D)$, with
uniformly positive, bounded couplings.  If $\mathcal A$ is any
increasing translation-invariant event on $\mathbb T_\ell$ and
$f_\ell(t)=\E\phi_{\mathbb T_\ell,t,\omega}(\mathcal A)$, then
\begin{equation}\label{crit:threshold-inequality}
 f_\ell'(t)\ge c\log\ell\, f_\ell(t)(1-f_\ell(t)),
 \qquad \ell\ge3,
\end{equation}
with $c>0$ independent of $\ell$ and $\mathcal A$.
\end{lemma}

\begin{proof}
For a fixed environment abbreviate its FK law to $\phi$, and set
$X_e=\one_{\{e\in A\}}$ and
$c_e(\omega)=\Cov_{\phi_{t,\omega}}(\one_{\mathcal A},X_e)\ge0$.
Differentiating the finite sum in \eqref{crit:fk-law} gives
\begin{equation}\label{crit:derivative}
 f_\ell'(t)=\sum_e\E[\lambda_e(t,\omega)c_e(\omega)],
 \qquad
 \lambda_e=\frac{\dd}{\dd t}\log v_e.
\end{equation}
Since $\dd K_e/\dd t=\tfrac12\tanh(2K_e)$, the coefficients
$\lambda_e=\tanh(2K_e)e^{2K_e}/(e^{2K_e}-1)$ have uniform positive
lower and upper bounds.

Represent each environment bit by a uniform variable $V_e$, with
$\xi_e=1$ when $V_e\ge1/2$.  Order the FK edges arbitrarily as
$e_1,\ldots,e_M$ and use further independent uniforms $U_i$ to set
\[
 X_{e_i}=\one_{\{U_i\ge1-p_i\}},\qquad
 p_i=\phi_{t,\omega}(X_{e_i}=1\mid X_{e_1},\ldots,X_{e_{i-1}}).
\]
The conditional FK laws are increasing both in the prescribed past
and in the couplings.  Induction in $i$ therefore makes the entire
sample increasing in all the $U_i$ and $V_e$.  Let $F$ be the Boolean
indicator of $\mathcal A$ in this product representation.  For one
uniform coordinate $z$, its influence $I_z$ is the probability, over
all other coordinates, that the section of $F$ in $z$ has values zero
and one on sets of positive Lebesgue measure.

We first bound an FK-uniform influence.  Fix the environment and let
$P=(X_{e_1},\ldots,X_{e_{i-1}})$ be the past.  Write
\[
 p(P)=\phi(X_{e_i}=1\mid P),\qquad
 a_b(P)=\phi(\mathcal A\mid P,X_{e_i}=b),\quad b=0,1.
\]
The future uniforms couple the two branches monotonically.  As both
intervals for $X_{e_i}$ have length at least $\eta$, the section in
$U_i$ is nonconstant exactly when the two resulting Boolean outputs
differ.  Averaging those future uniforms gives $a_1(P)-a_0(P)$.
Consequently
\begin{align}
 I_{U_i}
 &=\E\,\phi\bigl(a_1(P)-a_0(P)\bigr)\nonumber\\
 &\le\frac{1}{\eta(1-\eta)}
       \E\,\phi\bigl(\Cov_\phi(\one_{\mathcal A},X_{e_i}\mid P)\bigr)
 \le C\E c_{e_i}(\omega).\label{crit:uniform-influence}
\end{align}
For the last inequality, total covariance gives an additional term
\[
 \Cov_\phi\bigl(\phi(\mathcal A\mid P),
                        \phi(X_{e_i}=1\mid P)\bigr)\ge0.
\]
Both functions are increasing in $P$, and the marginal distribution
of $P$ is positively associated.  This proves the indicated sign.

For the environment coordinate $V_e$, condition on every other bond.
Let $\phi_-$ and $\phi_+$ be the FK laws with its low and high value,
and put $R=v_e^+/v_e^-\ge1$.  Changing this bond multiplies the weight
by $R^{X_e}$, so direct normalization gives
\begin{equation}\label{crit:bond-tilt}
 \phi_+(\mathcal A)-\phi_-(\mathcal A)
 =\frac{(R-1)\Cov_{\phi_-}(\one_{\mathcal A},X_e)}
              {1+(R-1)\phi_-(X_e)}.
\end{equation}
Monotonicity of the common-uniform coupling makes the left side the
conditional influence of $V_e$.  The ratio $R$ is uniformly bounded.
After averaging the other bonds, the low-bond covariance is at most
twice $\E c_e(\omega)$, since both low- and high-bond covariances are
nonnegative.  Thus $I_{V_e}\le C\E c_e(\omega)$ as well.

Translation invariance of the averaged law and of $\mathcal A$ makes
$\E c_e(\omega)$ constant within each of the two edge orientations.
Each orientation has $\ell^2$ edges.  Equation~\eqref{crit:derivative}
therefore implies
\begin{equation}\label{crit:max-influence}
 \max_z I_z\le C\ell^{-2}f_\ell'(t).
\end{equation}
There is no symmetry requirement on the chosen ordering or the
sequential sampler: only the averaged covariances were compared by
translation.  There are $N=2M=4\ell^2$ product coordinates.  The BKKKL
inequality~\cite{BKKKL1992}, in the null-set-invariant formulation
\cite[Theorem~2.3 and Equation~(2.3)]{GJN2013}, gives
\[
 \max_z I_z\ge c f_\ell(t)(1-f_\ell(t))\frac{\log N}{N}.
\]
Combine this with \eqref{crit:max-influence} to obtain
\eqref{crit:threshold-inequality}.
\end{proof}

\begin{proposition}\label{crit:tuning-zero}
For every sufficiently small fixed $D>0$, the tuned parameter is
$t_*(D)=0$.
\end{proposition}

\begin{proof}
The dual transformation $v_{e^*}^*=2/v_e$ is equivalent to
$\sinh(2K_e)\sinh(2K_{e^*}^*)=1$.  Thus, in the coordinates
\eqref{eq:model-td}, it sends $(t,D,\xi)$ to $(-t,D,-\xi)$ on the dual
lattice.  At $t=0$ the environment law is invariant.  Write
$a_i=\E\phi_{\mathbb T_\ell,0,\omega}(r(A)=i)$.
Lemma~\ref{crit:torus-duality} gives
$a_0/4\le a_2\le4a_0$.  In particular,
$a_1+a_2\ge1/5$ and $a_0+a_1\ge1/5$.

There are three nonzero characters of $\mathbb F_2^2$.  A one-dimensional
homology space is detected by two of them, whereas a two-dimensional
space is detected by all three.  Therefore
\begin{align*}
 \sum_{\chi\ne0}\E\phi_{0,\omega}(\mathcal H_\chi)
      &=2a_1+3a_2,\\
 \sum_{\chi\ne0}\E\phi_{0,\omega}(\mathcal H_\chi^c)
      &=3a_0+a_1.
\end{align*}
It follows that at $t=0$ at least one character has probability bounded
below by an absolute positive constant, and at least one has
probability bounded away from one.  The two characters need not be
the same.

Apply Lemma~\ref{crit:threshold} to each $\mathcal H_\chi$ along
$\ell_J=Qr_J$.  On $(0,1)$ it states
\[
 \frac{\dd}{\dd t}\log\frac{f_{\ell_J}(t)}{1-f_{\ell_J}(t)}
                         \ge c\log\ell_J.
\]
If $t_*>0$, integrate from $0$ to $t_*$ and use the upper bound in
\eqref{crit:tuned-homology}.  Every character probability at zero
would tend to zero.  If $t_*<0$, integrate from $t_*$ to zero and use
the lower bound in \eqref{crit:tuned-homology}; every character
probability at zero would tend to one.  Both conclusions contradict
the preceding absolute bounds.  Hence $t_*=0$.
\end{proof}

\subsection{Arm decay below zero and percolation above zero}

Proposition~\ref{read:arm} and Proposition~\ref{crit:tuning-zero} already
imply zero plus magnetization at $t=0$.  The next argument extracts a
stronger consequence strictly below zero.  It applies the monotone
OSSS inequality in each good environment, so it does not require a
sharpness theorem for a translation-invariant deterministic coupling
field.

\begin{lemma}[Averaged subcritical arm bound]\label{crit:subcritical-arm}
For every fixed $t'<0$ in the parameter window and sufficiently small
fixed $D>0$,
\begin{equation}\label{crit:subcritical-decay}
 \E\phi^{\mathrm w}_{t',\omega}
                  (0\longleftrightarrow\partial\Lambda_n)
                      \le C_{t',D}n^{-6}.
\end{equation}
\end{lemma}

\begin{proof}
Put $m=\lfloor n^{1/2}\rfloor$ and, for $z\in\Lambda_{2n}$, define
\[
 q_z(\omega)=\phi^{\mathrm w}_{z+\Lambda_m,0,\omega}
           (z\longleftrightarrow z+\partial\Lambda_m).
\]
Use the wall convention of Proposition~\ref{read:arm}; shifting the wall
one lattice step does not affect its uniform estimate.  Let
$\mathcal G_n=\{q_z\le n^{-1/50}\text{ for every }z\in\Lambda_{2n}\}$.
Markov's inequality, the $400$th moment in Proposition~\ref{read:arm},
and a union bound give
\begin{align}\label{crit:exception}
 \Prob(\mathcal G_n^c)
 &\le Cn^2 n^{400/50}m^{-400/8+o(1)}
   =n^{-15+o(1)}\le Cn^{-6}.
\end{align}
All boxes use the same global environment; independence of these
box events is unnecessary.  On $\mathcal G_n$, the same inequalities
hold throughout $t\in[t',0]$ by monotonicity in $t$.

Fix such an environment, and consider
\[
 f_\omega(t)=\phi^{\mathrm w}_{\Lambda_{4n},t,\omega}(\mathcal A_n),
 \qquad
 \mathcal A_n=\{0\longleftrightarrow\partial\Lambda_n
                        \text{ by a path in }\Lambda_n\}.
\]
Choose $k$ uniformly from $\{1,\ldots,n\}$, independently of the FK
configuration.  Explore all open components in $\Lambda_n$ meeting
its square shell $\partial\Lambda_k$, revealing all edges of the
induced graph on $\Lambda_n$ incident to the vertices reached.
Retain their actual connectivity; the shell
vertices are not identified with one another.  Every path witnessing
$\mathcal A_n$ meets the chosen shell, and its whole open component
is explored.  Thus this algorithm determines $\mathcal A_n$.

For an edge $e=xy$ within $\Lambda_n$, the chance that the chosen shell
lies within distance $m+2$ of $x$ or $y$ is at most $Cm/n$.  Otherwise,
querying $e$ requires that one endpoint is already joined to the shell
by a revealed open path, and hence has an open arm of radius $m$.
The endpoint box is contained in $\Lambda_{4n}$.  Conditional on all
edges outside that box, the domain Markov property gives an induced
boundary partition, dominated by full wiring.  Its conditional arm
probability is consequently at most $q_x$ or $q_y$, uniformly for
$t\in[t',0]$.  The revealment $\delta_e(t)$ therefore satisfies
\begin{equation}\label{crit:revealment}
 \max_e\delta_e(t)
          \le C(n^{-1/2}+n^{-1/50})\le Cn^{-1/50}.
\end{equation}
Edges outside $\Lambda_n$ are never queried.

The monotone-measure OSSS inequality
\cite[Theorem~1.1]{DCRT2019}, averaged over the independent choice of
shell, gives
\[
 f_\omega(t)(1-f_\omega(t))
 \le\sum_e\delta_e(t)
       \Cov_{\phi^{\mathrm w}_{\Lambda_{4n},t,\omega}}
                         (\one_{\mathcal A_n},X_e).
\]
The shell construction is the exploration underlying
\cite[Lemma~3.2]{DCRT2019}; here its revealment is bounded using the
quenched local arms on $\mathcal G_n$.  Differentiating the FK weights
as in \eqref{crit:derivative} and using their positive lower derivative
bound yields
\[
 f_\omega'(t)\ge c n^{1/50}f_\omega(t)(1-f_\omega(t)).
\]
The event that all four edges incident to zero are closed has
probability at least $\eta^4$, by sequential conditioning and finite
energy.  Hence $f_\omega(0)\le1-\eta^4$.  Integration of the preceding
logit inequality gives
\[
 f_\omega(t')\le C\exp\{-c|t'|n^{1/50}\}.
\]
After averaging, \eqref{crit:exception} contributes at most $Cn^{-6}$.
Finally, the restriction of the infinite wired law to this cylinder
event is dominated by the wired law in $\Lambda_{4n}$, by domain
comparison.  This proves \eqref{crit:subcritical-decay}.
\end{proof}

\begin{proposition}[Transition in the $(t,D)$ coordinates]
\label{crit:identification}
For every sufficiently small fixed $D>0$, the disorder-averaged plus
magnetization vanishes at $t\le0$ and is positive at every $t>0$.
\end{proposition}

\begin{proof}
At zero, Proposition~\ref{read:arm} implies
$\E\phi^{\mathrm w}_{\Lambda_n,0,\omega}
 (0\longleftrightarrow\partial\Lambda_n)\to0$ by H\"older's
inequality.  Equations~\eqref{crit:magnetization} and monotonicity give
zero magnetization for $t\le0$.

Fix $s>0$ small enough that $-s$ lies in the parameter window.  Planar
free/wired duality sends the infinite free FK law at $(s,D,\xi)$ to
the infinite wired dual law at $(-s,D,-\xi)$.  This identity follows
first from the finite planar bond expansion and then by taking free
and wired limits; see also the boundary conventions in
\cite[Section~2.5]{CDH2016}.  The dual environment still has the same
independent binary law.  Failure of a primal crossing of a
fixed-aspect rectangle of diameter comparable to $n$ entails a dual
crossing in the transverse direction, and hence a dual arm of length
at least $cn$ from a vertex in that rectangle.  There are $O(n^2)$
possible vertices.  Lemma~\ref{crit:subcritical-arm}, translated to
each such vertex and applied at $-s$, therefore bounds the averaged
failure probability by $Cn^{-4}$.

For each $n=2^kn_0$, use the four rectangular sides of
$\Lambda_{2n}\setminus\Lambda_n$, each of width $n$ and length $4n$.
Their longitudinal crossings intersect in the corner squares and
contain an open circuit surrounding $\Lambda_n$.  Require also a
left--right crossing of
$[n/2,5n]\times[-n/4,n/4]$.  Its initial point lies inside
$\Lambda_n$ and its endpoint lies outside $\Lambda_{4n}$, so it must
intersect both this circuit and the circuit at scale $2n$.  Rounding
rectangle endpoints by a lattice step is harmless; one may choose
$n_0$ divisible by four.  All aspect ratios are fixed and there are
five requested crossings per scale.  The probability that any of
them fails is at most
\[
 C\sum_{k\ge0}(2^kn_0)^{-4},
\]
which is smaller than one for large $n_0$.  On the complementary event
the joined circuits form an infinite open connected set.  Thus the
joint law obtained by first sampling the environment and then its
infinite free FK configuration has positive probability of an infinite
cluster.  This joint law is translation stationary.  If the origin
had zero probability of belonging to an infinite cluster, every
vertex would have probability zero, and countable subadditivity would
exclude an infinite cluster anywhere.  Hence the origin has positive
connection probability.  Wired domination and
\eqref{crit:magnetization} give positive averaged plus magnetization
at $s$.  Monotonicity extends this conclusion to every $t>0$.
\end{proof}

\subsection{The physical inverse temperature and the correlation ratio}

We finish by matching the parameter just identified to the original
bonds $1+\eps\xi_e$.  For $\eps\in(0,1)$ define
\[
 b_\pm(\beta,\eps)
       =\log\sinh\bigl(2\beta(1\pm\eps)\bigr).
\]
The mean $(b_++b_-)/2$ is continuous and strictly increasing from
$-\infty$ to $+\infty$ as $\beta$ runs through $(0,\infty)$.  It has
a unique zero $\beta_*(\eps)$, characterized by
\begin{equation}\label{crit:physical-equation}
 \sinh\bigl(2\beta_*(\eps)(1+\eps)\bigr)
 \sinh\bigl(2\beta_*(\eps)(1-\eps)\bigr)=1.
\end{equation}
Continuity and strict monotonicity imply
$\beta_*(\eps)\to\beta_0$ as $\eps\to0$.  At this parameter
put
\[
 D_*(\eps)=\tfrac12
       \bigl(b_+(\beta_*,\eps)-b_-(\beta_*,\eps)\bigr).
\]
Then $D_*(\eps)>0$ for $\eps>0$ and
$D_*(\eps)\to0$.  Choose $\eps_0>0$ so that this value belongs
to the permitted disorder interval whenever $0<\eps<\eps_0$.
At $\beta_*$ the physical model is exactly $(t,D)=(0,D_*)$.
Proposition~\ref{crit:identification} makes its magnetization zero;
monotonicity gives the same conclusion below $\beta_*$.

If $\beta>\beta_*$, both increments
$b_\pm(\beta,\eps)-b_\pm(\beta_*,\eps)$ are strictly positive.
Choose $s>0$, inside the parameter window, smaller than their minimum.
With the same environment, the physical couplings at $\beta$ dominate
those with log-sinh values $s+D_*\xi_e$.  The latter model has positive
averaged magnetization by Proposition~\ref{crit:identification}.
Consequently
\begin{equation}\label{crit:physical-beta}
 \beta_c(\eps)=\beta_*(\eps),
 \qquad m_\eps(\beta_c(\eps))=0.
\end{equation}
This comparison keeps $D_*$ fixed; no assumption about holding the
physical log-sinh difference constant while varying $\beta$ is used.

\begin{proof}[Proof of Theorem~\ref{thm:main}]
Fix $\eps\in(0,\eps_0)$ and use $D=D_*(\eps)>0$.
By \eqref{crit:physical-beta} and Proposition~\ref{crit:tuning-zero},
the tuned free infinite-volume correlation of
Proposition~\ref{read:moments} is $C_{\eps,\omega}(r)$ at the physical
critical inverse temperature.  That proposition gives, at its stopping
scale $J=\log r/\log L+O(1)$,
\[
 r_J^{2d}\E C_{\eps,\omega}(r)\asymp1,
 \qquad
 r_J^{4d}\E C_{\eps,\omega}(r)^2=J^{1/4+o(1)}.
\]
The factors $r_J^{4d}$ cancel in the ratio of the second moment to the
square of the first.  Thus
\[
 \log R_\eps(r)=(1/4+o(1))\log J+O(1),
 \qquad
 \frac{\log J}{\log\log r}\longrightarrow1,
\]
which proves the asserted limit along all integer $r\to\infty$.
The disorder remains fixed throughout this final limit.  The
thermodynamic limits and the separate Gibbs normalizations required
in the definition of $R_\eps$ are those already established in
Proposition~\ref{read:moments}.
\end{proof}
